\clearpage\SourceLeaf{450}%
\begin{center}\begingroup\fontsize{8}{9}\selectfont
\setlength{\tabcolsep}{4pt}
\begin{tabular}{r|r|p{170pt}|r}\hline
\(+\) südlich&\(+\) westlich&&Nr.\\\hline
\(-242773.3\)&\(+36827.3\)&Drochtersen&\((1)\)\\
\(-235431.6\)&\(+29726.7\)&Colmar&\((1)\)\\
\(-245551.2\)&\(+90763.0\)&Cappeln&\((1)\)\\
\(-245775.5\)&\(+87628.9\)&Midlum&\((1)\)\\
\(-245821.6\)&\(+24600.1\)&Neuendorf&\((1)\)\\
\(-246532.7\)&\(+67702.1\)&Oster Ilienworth&\((1)\)\\
\(-246892.5\)&\(+76614.2\)&Wanna&\((1)\)\\
\(-247388.5\)&\(+63676.7\)&Bülkau&\((1)\)\\
\(-247703.2\)&\(+36752.9\)&Krautsand&\((1)\)\\
\(-250722.6\)&\(+69097.6\)&Neuenkirchen&\((1)\)\\
\(-250784.2\)&\(+65963.9\)&Osterbruch&\((1)\)\\
\(-251343.0\)&\(+34554.1\)&Glückstadt kleiner Turm (dänische Station)&\((1)\)\\
\(-251399.8\)&\(+30139.2\)&Herzhorn&\((1)\)\\
\(-251523.1\)&\(+11052.9\)&Barmstedt&\((1)\)\\
\(-251573.4\)&\(+34126.7\)&Glückstadt Kirchturm&\((1)\)\\
\(-252146.1\)&\(+40748.4\)&Hammelvörden&\((1)\)\\
\(-253330.1\)&\(+78043.0\)&Lüdingworth&\((1)\)\\
\(-254110.4\)&\(+21378.6\)&Horst&\((1)\)\\
\(-254115.7\)&\(+68778.6\)&Otterndorf&\((1)\)\\
\(-255131.8\)&\(+28321.4\)&Süderau&\((1)\)\\
\(-255141.2\)&\(+62459.0\)&Belum&\((1)\)\\
\(-255213.6\)&\(+83123.0\)&Franzenburg&\((1)\)\\
\(-255856.6\)&\(+43050.3\)&Freiburg&\((1)\)\\
\hline\end{tabular}\endgroup\end{center}
\clearpage
\begin{center}\begingroup\fontsize{8}{9}\selectfont
\setlength{\tabcolsep}{4pt}
\begin{tabular}{r|r|p{170pt}|r}\hline
\(+\) südlich&\(+\) westlich&&Nr.\\\hline
\(-255908.9\)&\(+33800.4\)&Borsfleth&\((1)\)\\
\(-256062.6\)&\(+84523.8\)&Hartmanns Platz 1 bei Altenwalde&\((1)\)\\
\(-256062.7\)&\(+84104.8\)&Altenwalde&\((1)\)\\
\(-256069.1\)&\(+77120.7\)&Altenbruch Spitze 1.&\((1)\)\\
\(-256075.2\)&\(+77119.1\)&Altenbruch Spitze 2.&\((1)\)\\
\(-256894.2\)&\(+29816.6\)&Crempe&\((1)\)\\
\(-257722.7\)&\(+21032.0\)&Hohenfelde&\((1)\)\\
\(-257778.1\)&\(+35721.5\)&Wevelsfleth&\((1)\)\\
\(-257996.5\)&\(+80031.2\)&Groden&\((1)\)\\
\(-258998.8\)&\(+15610.9\)&Hörnerkirchen&\((1)\)\\
\(-259523.9\)&\(+26920.1\)&Neuenbrook&\((1)\)\\
\(-259779.2\)&\(+40370.0\)&Brockdorf&\((1)\)\\
\(-259854.2\)&\(+81878.1\)&Ritzebüttel Giebelstange 1&\((1)\)\\
\(-260086.9\)&\(+33836.4\)&Neuenkirchen&\((1)\)\\
\(-261494.4\)&\(+81196.7\)&Cuxhaven Leuchtturm&\((1)\)\\
\(-261725.0\)&\(+34726.2\)&Beienfleth&\((1)\)\\
\(-262754.7\)&\(+83558.8\)&Döse&\((1)\)\\
\(-263188.1\)&\(+45161.9\)&St. Margareth&\((1)\)\\
\(-266552.1\)&\(+37341\,8\)&Wilster&\((1)\)\\
\(-266569.396\)&\(+95074.242\)&Neuwerk Leuchtturm Zentrum&\((1)\)\\
\(-269380.6\)&\(+14702.3\)&Kellinghusen&\((1)\)\\
\(-270223.4\)&\(+61150.8\)&Marne&\((1)\)\\
\hline\end{tabular}\endgroup\end{center}

Zur Bedeutung der Koordinaten ist Folgendes zu bemerken.

Will man sich nur eine allgemeine Vorstellung davon machen, so kann man die erste Zahl als Angabe dafür ansehen, wie weit der betreffende Ort südlich (beim \(+\)-Zeichen) oder nördlich (beim \(-\)-Zeichen) der Göttinger Sternwarte liegt. Die zweite Zahl gibt dagegen an, wie weit er westlich (bei \(+\)) oder östlich (bei \(-\)) liegt.

Die Krümmung der Erdoberfläche ist dabei bereits so berücksichtigt, dass sich beim Eintragen dieser Koordinaten in eine Ebene ein \textit{konformes}, das heißt in den kleinsten Teilen ähnliches Bild ergibt. Meine geodätischen Abhandlungen enthalten bereits einen Teil der näheren Ausführungen dazu, und spätere Abhandlungen werden dies noch ausführlicher darlegen.

Der genaue Ursprung der Koordinaten in der Sternwarte ist übrigens der Mittelpunkt der Achse des Reichenbachschen Meridiankreises.

Als Einheit der Koordinaten wurde diejenige Länge gewählt, die nach dem besten Kenntnisstand von 1821 als der zehnmillionste Teil des Quadranten des Erdmeridians gelten konnte, nämlich 443,307885 Pariser Linien. Sie unterscheidet sich etwas, wenn auch nur sehr wenig, vom sogenannten gesetzlichen französischen Meter. Dieser war bekanntlich auf 443,296 Pariser Linien festgesetzt worden. Obwohl in späterer Zeit (seit 1821) weitere Bestimmungen des zehnmillionsten Teils des Erdmeridianquadranten gewonnen wurden, und zwar stets deutlich größer als der eben angeführte gesetzliche Meter), habe ich es doch vorgezogen, bei der einmal gewählten Einheit zu bleiben, da man jede einzelne Zahl leicht in jede beliebige andere Einheit umrechnen kann, sobald das Verhältnis bekannt ist.

